% GENERATED FILE. Edit canonical lessons, then run npm run generate:books.
% canonical-source-hash: fnv1a64:311f2745d8549a22
% canonical-lessons: BN-C397-ghat, BN-C397-lonch, BN-C397-kheya, BN-C397-par, BN-C397-bherano

\chapter{Ghats and Ferries}
\label{ch:bn-ghats-and-ferries}
\hlchaptermodality{397}

\begin{chapteropening}
\textbf{By the end of this chapter:} \emph{I can cross a river by ferry from one ghat to the other bank.}

Complete the last lesson of chapter 397: I can cross a river by ferry from one ghat to the other bank.
\end{chapteropening}

\section[\texorpdfstring{\bn{ঘাট} (ghāṭ)}{ঘাট (ghāṭ)}]{\bn{ঘাট} (ghāṭ) — a landing place; riverside steps; a ghat}

\label{lesson:BN-C397-ghat}



\begin{quote}
Before the new one: say the Bengali for the hood of a country boat, then the Bengali for to row.
\end{quote}

\subsection*{You'll want to know: \bn{ঘাট}}
\textbf{\bn{ঘাট}} — \emph{ghāṭ} — \textquotedblleft{}a landing place; riverside steps; a ghat\textquotedblright{}.

\begin{quote}
\bn{সকালে} \bn{ঘাটে} \bn{অনেক} \bn{লোক} \bn{স্নান} \bn{করে।}
\end{quote}

\begin{quote}
\emph{sôkāle ghāṭe ônek lok snān kôre.}
\end{quote}

\begin{quote}
In the morning many people bathe at the ghat.
\end{quote}

The steps down to a river or pond, and the place where boats come in.

\begin{cousinweb}
\textbf{\bn{ঘাট}} (\emph{ghāṭ}) is inherited: CDIAL lists Bengali \emph{ghāṭ}, \textquotedblleft{}a mountain pass, a landing place\textquotedblright{}, under Sanskrit \emph{ghaṭṭa}, \textquotedblleft{}a landing place, a quay\textquotedblright{} (CDIAL 4414); English \emph{ghat} is the same word, borrowed.
\end{cousinweb}

\begin{grammarlens}[title={What you've built}]
One more word for getting around.
\end{grammarlens}

\subsection*{Your turn}
\begin{itemize}
\raggedright
  \item \emph{Say it:} \emph{ghāṭ}
  \item \emph{Say it:} \emph{ghāṭ}, once more
  \item \emph{Say it:} \emph{sôkāle ghāṭe ônek lok snān kôre.} — In the morning many people bathe at the ghat.
  \item \emph{Recall:} say the Bengali for the hood of a country boat, then the Bengali for to row, then say \emph{ghāṭ} again
\end{itemize}

\begin{tcolorbox}[breakable,colback=teal!4,colframe=teal!35!black,title={Before you move on}]
What does \bn{ঘাট} mean? (\textquotedblleft{}A landing place; riverside steps; a ghat\textquotedblright{}.) Say it once more.
\end{tcolorbox}
\section[\texorpdfstring{\bn{লঞ্চ} (lônch)}{লঞ্চ (lônch)}]{\bn{লঞ্চ} (lônch) — a launch; a motor ferry}

\label{lesson:BN-C397-lonch}



\begin{quote}
Before the new one: say the Bengali for to row, then the Bengali for a landing place.
\end{quote}

\subsection*{You'll want to know: \bn{লঞ্চ}}
\textbf{\bn{লঞ্চ}} — \emph{lônch} — \textquotedblleft{}a launch; a motor ferry\textquotedblright{}.

\begin{quote}
\bn{লঞ্চটা} \bn{ভোরে} \bn{ঘাট} \bn{ছাড়ল।}
\end{quote}

\begin{quote}
\emph{lônchṭā bhore ghāṭ chhāṛlô.}
\end{quote}

\begin{quote}
The launch left the ghat at dawn.
\end{quote}

The motor boat that carries passengers on the big rivers.

\begin{cousinweb}
\textbf{\bn{লঞ্চ}} (\emph{lônch}) is borrowed from English \emph{launch}, a motor boat; \bn{ঞ্চ} is the n-ch joined letter.
\end{cousinweb}

\begin{grammarlens}[title={What you've built}]
One more word for getting around.
\end{grammarlens}

\subsection*{Your turn}
\begin{itemize}
\raggedright
  \item \emph{Say it:} \emph{lônch}
  \item \emph{Say it:} \emph{lônch}, once more
  \item \emph{Say it:} \emph{lônchṭā bhore ghāṭ chhāṛlô.} — The launch left the ghat at dawn.
  \item \emph{Recall:} say the Bengali for to row, then the Bengali for a landing place, then say \emph{lônch} again
\end{itemize}

\begin{tcolorbox}[breakable,colback=teal!4,colframe=teal!35!black,title={Before you move on}]
What does \bn{লঞ্চ} mean? (\textquotedblleft{}A launch; a motor ferry\textquotedblright{}.) Say it once more.
\end{tcolorbox}
\section[\texorpdfstring{\bn{খেয়া} (kheyā)}{খেয়া (kheyā)}]{\bn{খেয়া} (kheyā) — a ferry (across a river)}

\label{lesson:BN-C397-kheya}



\begin{quote}
Before the new one: say the Bengali for a landing place, then the Bengali for a launch.
\end{quote}

\subsection*{You'll want to know: \bn{খেয়া}}
\textbf{\bn{খেয়া}} — \emph{kheyā} — \textquotedblleft{}a ferry (across a river)\textquotedblright{}.

\begin{quote}
\bn{আমরা} \bn{খেয়ায়} \bn{নদী} \bn{পার} \bn{হলাম।}
\end{quote}

\begin{quote}
\emph{āmrā kheyāy nôdī pār hôlām.}
\end{quote}

\begin{quote}
We crossed the river by ferry.
\end{quote}

A \bn{খেয়া} is a short crossing back and forth; a \bn{লঞ্চ} makes longer runs.

\begin{cousinweb}
\textbf{\bn{খেয়া}} (\emph{kheyā}) is inherited: CDIAL lists Bengali \emph{kheyā}, \textquotedblleft{}a ferry\textquotedblright{}, under Sanskrit \emph{kṣēpa}, \textquotedblleft{}a sending, a throw\textquotedblright{} (CDIAL 3738); the step from sending to the boat that sends you across is the Bengali side's own.
\end{cousinweb}

\begin{grammarlens}[title={What you've built}]
One more word for getting around.
\end{grammarlens}

\subsection*{Your turn}
\begin{itemize}
\raggedright
  \item \emph{Say it:} \emph{kheyā}
  \item \emph{Say it:} \emph{kheyā}, once more
  \item \emph{Say it:} \emph{āmrā kheyāy nôdī pār hôlām.} — We crossed the river by ferry.
  \item \emph{Recall:} say the Bengali for a landing place, then the Bengali for a launch, then say \emph{kheyā} again
\end{itemize}

\begin{tcolorbox}[breakable,colback=teal!4,colframe=teal!35!black,title={Before you move on}]
What does \bn{খেয়া} mean? (\textquotedblleft{}A ferry (across a river)\textquotedblright{}.) Say it once more.
\end{tcolorbox}
\section[\texorpdfstring{\bn{পাড়} (pāṛ)}{পাড় (pāṛ)}]{\bn{পাড়} (pāṛ) — a bank (of a pond or river); a border (of a sari)}

\label{lesson:BN-C397-par}



\begin{quote}
Before the new one: say the Bengali for a launch, then the Bengali for a ferry.
\end{quote}

\subsection*{You'll want to know: \bn{পাড়}}
\textbf{\bn{পাড়}} — \emph{pāṛ} — \textquotedblleft{}a bank (of a pond or river); a border (of a sari)\textquotedblright{}.

\begin{quote}
\bn{পুকুরের} \bn{পাড়ে} \bn{একটা} \bn{আমগাছ} \bn{আছে।}
\end{quote}

\begin{quote}
\emph{pukurer pāṛe ekṭā āmgāchh āchhe.}
\end{quote}

\begin{quote}
There is a mango tree on the bank of the pond.
\end{quote}

Also the coloured border of a sari: \bn{লাল} \bn{পাড়ের} \bn{শাড়ি}, a sari with a red border.

\begin{cousinweb}
\textbf{\bn{পাড়}} (\emph{pāṛ}) probably belongs with the reconstructed \emph{pāḍi}, under which CDIAL lists Bengali \emph{pāṛi}, \textquotedblleft{}an edge, a border\textquotedblright{} (CDIAL 8041): a bank is the edge of the water, a sari's \bn{পাড়} the edge of the cloth.
\end{cousinweb}

\begin{grammarlens}[title={What you've built}]
One more word for getting around.
\end{grammarlens}

\subsection*{Your turn}
\begin{itemize}
\raggedright
  \item \emph{Say it:} \emph{pāṛ}
  \item \emph{Say it:} \emph{pāṛ}, once more
  \item \emph{Say it:} \emph{pukurer pāṛe ekṭā āmgāchh āchhe.} — There is a mango tree on the bank of the pond.
  \item \emph{Recall:} say the Bengali for a launch, then the Bengali for a ferry, then say \emph{pāṛ} again
\end{itemize}

\begin{tcolorbox}[breakable,colback=teal!4,colframe=teal!35!black,title={Before you move on}]
What does \bn{পাড়} mean? (\textquotedblleft{}A bank (of a pond or river); a border (of a sari)\textquotedblright{}.) Say it once more.
\end{tcolorbox}
\section[\texorpdfstring{\bn{ভেড়ানো} (bheṛāno)}{ভেড়ানো (bheṛāno)}]{\bn{ভেড়ানো} (bheṛāno) — to bring (a boat) alongside; to moor}

\label{lesson:BN-C397-bherano}



\begin{quote}
Before the new one: say the Bengali for a ferry, then the Bengali for a bank (of a pond or river).
\end{quote}

\subsection*{You'll want to know: \bn{ভেড়ানো}}
\textbf{\bn{ভেড়ানো}} — \emph{bheṛāno} — \textquotedblleft{}to bring (a boat) alongside; to moor\textquotedblright{}.

\begin{quote}
\bn{মাঝি} \bn{নৌকাটা} \bn{ঘাটে} \bn{ভেড়াল।}
\end{quote}

\begin{quote}
\emph{mājhi noukāṭā ghāṭe bheṛālô.}
\end{quote}

\begin{quote}
The boatman brought the boat in to the ghat.
\end{quote}

Past (he/she …): \bn{তিনি} \bn{নৌকা} \bn{ঘাটে} \bn{ভেড়ালেন}, \emph{tini noukā ghāṭe bheṛālen}.

Not \bn{ভেড়া}, a sheep: the same first letters, a different word.

\begin{cousinweb}
\textbf{\bn{ভেড়ানো}} (\emph{bheṛāno}) is inherited: CDIAL lists Bengali \emph{bhiṛāna}, \textquotedblleft{}to bring alongside\textquotedblright{}, under the reconstructed \emph{bhiṭ}, \textquotedblleft{}meet, throng\textquotedblright{} (CDIAL 9490), the same root as \bn{ভিড়}, a crowd.
\end{cousinweb}

\begin{grammarlens}[title={What you've built}]
One more word for getting around.
\end{grammarlens}

\subsection*{Your turn}
\begin{itemize}
\raggedright
  \item \emph{Say it:} \emph{bheṛāno}
  \item \emph{Say it:} \emph{bheṛāno}, once more
  \item \emph{Say it:} \emph{mājhi noukāṭā ghāṭe bheṛālô.} — The boatman brought the boat in to the ghat.
  \item \emph{Recall:} say the Bengali for a ferry, then the Bengali for a bank (of a pond or river), then say \emph{bheṛāno} again
\end{itemize}

\begin{tcolorbox}[breakable,colback=teal!4,colframe=teal!35!black,title={Before you move on}]
What does \bn{ভেড়ানো} mean? (\textquotedblleft{}To bring (a boat) alongside; to moor\textquotedblright{}.) Say it once more.
\end{tcolorbox}
